\theory{Class field theory}{How can all abelian extensions of a number system be described using arithmetic information already present in that system?}{Class field theory is the part of algebraic number theory that describes abelian Galois extensions of local and global fields. It translates field extensions into ideal classes, norms, units, and ideles, and it gives reciprocity laws that generalize quadratic reciprocity.}{Prime numbers, algebraic number theory, explicit reciprocity, Iwasawa theory, the Langlands program, arithmetic geometry, and the study of how primes split in extensions.}{Reciprocity maps ideals or ideles to abelian Galois groups, turning prime-splitting questions into calculations in the base field.}

\paragraph{The problem and its history}

A field is a number system in which addition, subtraction, multiplication, and division by nonzero numbers are possible. A field extension is a larger system containing the original one. For example, adjoining $\sqrt{2}$ to the rational numbers gives $\mathbb Q(\sqrt2)$. Galois theory studies the symmetries of an extension: a Galois automorphism rearranges the new numbers while leaving the old numbers fixed.

Class field theory asks for a complete description when these symmetries commute with one another, the \emph{abelian} case. The surprising idea is that the answer can be read from arithmetic objects in the original field: ideals, their classes, local units, and a global object called the idele class group. \cite{class_field_theory_ref1,class_field_theory_ref4}

The story began with Gauss's quadratic reciprocity law, which relates whether one prime is a square modulo another to the reversed question. The later generalization involved quadratic forms, genus theory, Kummer's work, ideals and completions, and cyclotomic extensions. Kronecker and Weber recognized explicit patterns for the rational numbers; Weber coined the phrase class field, and Hilbert formulated central problems. Takagi proved the existence theorem, while Artin developed the reciprocity law, with Chebotarev's density theorem supplying a crucial description of primes in extensions. \cite{class_field_theory_ref2,class_field_theory_ref3}

\end{historylist}
\section*{3. Theory}
\subsection*{Core theory}
\paragraph{Local fields, global fields, and the central dictionary}

A number field is a finite extension of $\mathbb Q$, such as $\mathbb Q(i)$ or $\mathbb Q(\sqrt5)$. A global function field is a finite extension of a rational function field over a finite field. These are the global fields. A local field is obtained by completing a global field at one of its places. Completion means that we add limits of Cauchy sequences, just as the real numbers complete the rationals. For a prime $p$, the field $\mathbb Q_p$ remembers very precise information about divisibility by $p$.

The maximal abelian extension $A$ of a field $K$ is the largest algebraic extension whose Galois symmetries commute. Its Galois group is usually infinite, but it has a natural topology: finite pieces of information determine open neighborhoods. It is therefore a profinite group, an inverse limit of finite groups.

Class field theory constructs an arithmetic topological group associated with $K$. For a local field with finite residue field, the key object is the multiplicative group $K^\times$. For a global field, the key object is the idele class group $C_K$. An idele is a compatible collection of local numbers, one for each place, with a finiteness condition; quotienting by the diagonal copy of $K^\times$ removes the descriptions that come from a single global number. \cite{class_field_theory_ref1,class_field_theory_ref5}

The dictionary has two directions:

{\small
\[
\begin{array}{c|c}
\text{Arithmetic information in }K & \text{Abelian extension information}\\
\hline
\text{an open finite-index subgroup of }C_K & \text{a finite abelian extension of }K\\
\text{a norm subgroup from }L & \text{the extension }L/K\\
\text{an ideal class in an unramified situation} & \text{a Frobenius automorphism}
\end{array}
\]
}

This is useful because arithmetic groups can often be computed or approximated, whereas listing all algebraic extensions directly would be impossible.

\paragraph{The Hilbert class field}

Let $F$ be a number field. Its ideals form a multiplicative system, but not every ideal is generated by one element. The ideal class group measures this failure: two ideals have the same class when their quotient differs by multiplication by a principal ideal. The class group is finite, and its size is the class number.

The Hilbert class field $K$ of $F$ is the maximal abelian extension of $F$ that is unramified at every finite and infinite place. The Hilbert class field theorem gives a canonical isomorphism
\[
 \operatorname{Gal}(K/F)\cong \operatorname{Cl}(F).
\]
Thus the number of ideal classes is exactly the number of symmetries in the largest everywhere-unramified abelian extension. An ideal $\mathfrak p$ not dividing the discriminant determines a Frobenius automorphism in the Galois group; under the isomorphism, this automorphism is represented by the ideal class of $\mathfrak p$. \cite{class_field_theory_ref1,class_field_theory_ref2}

For a simple example, the rational field has class group of size one and no nontrivial everywhere-unramified abelian extension. A field with a nontrivial class group has a genuinely larger Hilbert class field. Computing ideals and their classes therefore reveals an extension field that may be difficult to write down directly.

\paragraph{Artin reciprocity and the existence theorem}

For a finite abelian extension $L/F$, the idelic norm map sends an idele of $L$ to an idele of $F$. Class field theory states that the quotient
\[
 \theta_{L/F}:C_F/N_{L/F}(C_L)\longrightarrow \operatorname{Gal}(L/F)
\]
is a canonical isomorphism. The map is the Artin reciprocity map. In the unramified prime case, it sends the class of a prime to its Frobenius automorphism, so the abstract map connects arithmetic with the way that prime acts on the roots of a polynomial modulo that prime. \cite{class_field_theory_ref1,class_field_theory_ref4}

The existence theorem says that finite abelian extensions of $F$ correspond bijectively to closed subgroups of finite index in $C_F$. The subgroup attached to $L$ is the norm group $N_{L/F}(C_L)$. In one direction, the extension produces a subgroup. In the other, a suitable subgroup produces a unique class field. This is why the theory is called class field theory: a class or subgroup of the arithmetic object determines a field.

The formulation also has an infinite version. The Galois group of the maximal abelian extension is naturally the profinite completion of $C_F$. Equivalently, the abelianization of the absolute Galois group is determined by the arithmetic class group. This unifies all finite abelian extensions instead of studying them one at a time. \cite{class_field_theory_ref1,class_field_theory_ref5}

\paragraph{Local and global reciprocity}

Local class field theory begins with a local field $K$. It gives a reciprocity isomorphism from $K^\times$ to the Galois group of its maximal abelian extension, interpreted with the correct profinite topology. The norm subgroup of a finite abelian extension $L/K$ is the kernel of the corresponding map. Local units describe the part of the extension that is ramified, while the valuation records the unramified direction.

To pass from local to global theory, complete a number field at every place and apply the local maps. Their product is trivial on a global element embedded diagonally in all local fields. This is the global reciprocity law. It generalizes the product formula and reaches far beyond quadratic reciprocity: local contributions at all primes and at infinity fit together to produce one global statement. \cite{class_field_theory_ref1,class_field_theory_ref3}

There are two broad proof styles. Artin and Tate developed a cohomological method using group cohomology and class formations. It gives a powerful conceptual organization of norms, extensions, and obstructions. Later, Neukirch and others developed more explicit and algorithmic approaches that avoid some cohomology, making the reciprocity maps easier to calculate in examples. The two approaches prove the same central results but illuminate different parts of the subject. \cite{class_field_theory_ref3,class_field_theory_ref6}

\paragraph{Explicit class field theory}

General class field theory tells us that extensions exist and what arithmetic group controls them. Explicit class field theory tries to write the extensions down.

The Kronecker--Weber theorem says that every finite abelian extension of $\mathbb Q$ is contained in a cyclotomic field $\mathbb Q(\zeta_n)$, where $\zeta_n$ is an $n$th root of unity. Thus roots of unity generate all abelian extensions of the rational numbers. For example, the quadratic field $\mathbb Q(i)$ lies in $\mathbb Q(\zeta_4)$, and the real quadratic field $\mathbb Q(\sqrt5)$ lies in the maximal real subfield of $\mathbb Q(\zeta_5)$. \cite{class_field_theory_ref1,class_field_theory_ref4}

For imaginary quadratic fields, elliptic curves with complex multiplication give another explicit source. Torsion points of these curves and special values of modular functions generate abelian extensions. Shimura's theory extends this explicit viewpoint to broader classes of fields. In positive characteristic, cyclotomic and Drinfeld-module constructions provide parallel descriptions, and Witt duality gives a particularly direct description of the $p$-part of reciprocity in some settings. \cite{class_field_theory_ref4}

These explicit constructions are powerful but special. Roots of unity are naturally available over $\mathbb Q$, and complex multiplication supplies extra structure for imaginary quadratic fields. A general number field usually lacks such a visible source of generators, which is why the idelic formulation is needed.

\paragraph{Applications, boundaries, and generalizations}

Class field theory explains how primes split, remain inert, or ramify in abelian extensions. The Frobenius map turns this into arithmetic data that can be used to test congruences and factorization patterns. The theory is also used in Iwasawa theory, Galois module theory, Artin--Verdier duality, and explicit studies of special values of $L$-functions. Many advances toward the Langlands correspondence, the Birch and Swinnerton-Dyer conjecture, and Iwasawa theory use explicit class field methods or their extensions. \cite{class_field_theory_ref1,class_field_theory_ref7}

Class field theory deliberately treats the abelian case. Nonabelian Galois groups do not fit into the same simple correspondence with subgroups of an abelian class group. The Langlands program is often viewed as a proposed nonabelian class field theory, but it does not currently provide the same existence theorem for class fields. Anabelian geometry asks whether a number field or a hyperbolic curve can be reconstructed from its full absolute Galois group or algebraic fundamental group. Higher class field theory studies abelian extensions of higher local and global fields; its arithmetic groups are higher $K$-groups such as Milnor $K$-groups rather than only $K_1=K^\times$. \cite{class_field_theory_ref1,class_field_theory_ref8}

Class field theory depends on field theory, Galois theory, ideals, prime factorization, completions, topological groups, and group cohomology. It helps algebraic number theory by providing a complete abelian extension dictionary. It also supplies the abelian base case for the Langlands program, local tools for arithmetic geometry, and explicit reciprocity laws used in the study of elliptic curves and $L$-functions.

\section*{4. Important theorems and results}
\begin{enumerate}[leftmargin=*,itemsep=0.35em]

\item \textbf{Hilbert class field theorem.} For a number field $F$, the Galois group of its maximal everywhere-unramified abelian extension is canonically isomorphic to the ideal class group of $F$. \cite{class_field_theory_ref1}

\item \textbf{Artin reciprocity law.} For every finite abelian extension $L/F$, the quotient $C_F/N_{L/F}(C_L)$ is canonically isomorphic to $\operatorname{Gal}(L/F)$. \cite{class_field_theory_ref1,class_field_theory_ref3}

\item \textbf{Existence theorem.} Finite abelian extensions of $F$ correspond to closed finite-index subgroups of the idele class group $C_F$. \cite{class_field_theory_ref1}

\item \textbf{Kronecker--Weber theorem.} Every finite abelian extension of $\mathbb Q$ is contained in a cyclotomic field. \cite{class_field_theory_ref4}

\item \textbf{Chebotarev density theorem.} In a finite Galois extension, prime ideals are distributed among Frobenius conjugacy classes with densities determined by the sizes of those classes. In the abelian case, this turns reciprocity classes into predictable families of primes. \cite{class_field_theory_ref2}
\end{enumerate}

\paragraph{Applications}
\begin{enumerate}[leftmargin=*,itemsep=0.35em]
\item \textbf{Cryptographic protocols using class groups.} Class groups of imaginary quadratic fields provide abelian groups of unknown structure, which is a desirable foundation for cryptographic constructions. The group $\operatorname{Cl}(\mathbb Q(\sqrt{-d}))$ is hard to compute explicitly for large square-free $d$, yet its order and structure can be verified using class field theory's prediction that the Hilbert class field has Galois group isomorphic to $\operatorname{Cl}(\mathbb Q(\sqrt{-d}))$. Cryptosystems based on this hardness include commitment schemes, verifiable delay functions, and group-action-based protocols analogous to Diffie--Hellman: two parties publicly exchange ideal classes and combine them with their secret ideals, arriving at a shared class without revealing the underlying ideals. The security rests on the assumption that computing the class number or finding a principal ideal representation is computationally intractable for appropriately chosen discriminants \cite{class_field_theory_ref1,class_field_theory_ref4}.
\item \textbf{Primality proving via class field theory.} The theory of complex multiplication gives an explicit construction of abelian extensions of imaginary quadratic fields, and the resulting class polynomials can be used to prove that an integer is prime. Specifically, one computes the Hilbert class polynomial $H_d(X)$ for discriminant $-d$ and verifies that a given elliptic curve $E$ over $\mathbb F_p$ has $j$-invariant equal to a root of $H_d(X)$ modulo $p$. Class field theory guarantees that $p$ splits completely in the Hilbert class field $\mathbb Q(\sqrt{-d},E)$ only when $p$ is prime and satisfies the splitting conditions predicted by Artin reciprocity. This method, implemented in software such as APRCL and CM-based primality tests, proves primality of numbers with thousands of digits and is part of the infrastructure used by cryptographic libraries to generate large prime moduli for RSA and Diffie--Hellman key exchange \cite{class_field_theory_ref2,class_field_theory_ref3}.
\item \textbf{Ideal class group and arithmetic dynamics.} The growth of denominators in orbits of rational points under polynomial iteration is governed by height functions whose local components are determined by splitting behavior in abelian extensions. Class field theory predicts how primes ramify or remain inert in these extensions, which in turn controls the canonical height that measures orbit growth. This connection has been used to prove that certain polynomial maps have only preperiodic rational points---a dynamical analogue of the Mordell conjecture---by reducing the height computation to Artin reciprocity and Chebotarev density arguments \cite{class_field_theory_ref1,class_field_theory_ref7}.
\end{enumerate}
\theoryconnections{class field theory}{%
\item Galois theory supplies extension groups, algebraic number theory supplies ideals and units, and local fields supply completions at primes and infinite places.
\item Topological groups and harmonic analysis are needed because the relevant global Galois groups are profinite and the idele class group is topological.
}{%
\item Class field theory helps arithmetic geometry and Diophantine analysis by predicting how primes split in extensions.
\item Its reciprocity map also provides the abelian part of the wider Langlands viewpoint, where arithmetic information is related to representations and analytic functions.
}

\section*{Glossary}
\begin{description}[leftmargin=*,style=nextline,itemsep=0.3em]

\item[\textbf{Ideal class group.}] The group that measures the failure of unique factorization into principal ideals in a number field; its size (the class number) counts how many distinct classes of ideals exist.

\item[\textbf{Abelian extension.}] A field extension whose Galois symmetries commute with one another; class field theory provides a complete description of all such extensions from arithmetic data in the base field.

\item[\textbf{Reciprocity map.}] A canonical isomorphism from an arithmetic group (such as the idele class group) to the Galois group of an abelian extension, translating prime-splitting behavior into abstract symmetry.

\item[\textbf{Idele class group.}] The quotient of the idele group of a number field by its diagonal copy of the field's nonzero elements; it is the arithmetic object that class field theory pairs with the abelian Galois group.

\item[\textbf{Frobenius automorphism.}] The Galois automorphism associated to an unramified prime in an extension, which raises residue field elements to the power of the prime's cardinality; it records how that prime acts on the extension.

\item[\textbf{Hilbert class field.}] The maximal abelian extension of a number field that is unramified at every place; its Galois group is isomorphic to the ideal class group of the base field.

\item[\textbf{Ramification.}] The phenomenon in which a prime ideal of the base field fails to remain prime or to have the expected number of factors in an extension; ramified primes are the ones where the extension ``wraps around.''

\item[\textbf{Galois group.}] The group of field automorphisms of an extension that fix the base field; class field theory studies its abelian quotient.

\item[\textbf{Number field.}] A finite extension of $\mathbb Q$, such as $\mathbb Q(i)$; class field theory describes abelian extensions of number fields.

\item[\textbf{Completion.}] Adding limits of Cauchy sequences at a given place, producing a local field; bridges global and local class field theory.

\item[\textbf{Profinite group.}] A topological group that is an inverse limit of finite groups, hence compact, Hausdorff, and totally disconnected.

\item[\textbf{Chebotarev density theorem.}] In a finite Galois extension, Frobenius elements are distributed among conjugacy classes with densities proportional to class sizes.

\end{description}

\section*{7. Sample Questions}
\emph{Exercise reference:} \cite{class_field_theory_ref1}. The cited edition is the source for the exercise set; consult its Exercises section for the official statement and numbering.
\begin{enumerate}[leftmargin=*,itemsep=0.9em]

\item \textbf{Exercise 1.} Compute the ideal class group of $F=\mathbb{Q}(\sqrt{-5})$. Specifically, show that the ideal $(2,1+\sqrt{-5})$ is non-principal.

\emph{Solution.} The ring of integers is $\mathbb{Z}[\sqrt{-5}]$. The ideal $I=(2,1+\sqrt{-5})$ has norm $N(I)=2$ (since $2$ ramifies: $(2)=(2,1+\sqrt{-5})^2$). If $I$ were principal, say $I=(\alpha)$, then $N(\alpha)=N(I)=2$, so $a^2+5b^2=2$ for $\alpha=a+b\sqrt{-5}$. This has no integer solutions (since $a^2+5b^2\ge 5$ for $|b|\ge 1$ and $a^2\le 2$ gives $a=0,\pm 1$, none yielding $a^2+5b^2=2$). So $I$ is non-principal, and the class number is at least 2. In fact $h(\mathbb{Q}(\sqrt{-5}))=2$.

\item \textbf{Exercise 2.} Use the Kronecker--Weber theorem to show that $\mathbb{Q}(\sqrt{5})$ is contained in a cyclotomic field. Find the smallest such cyclotomic field.

\emph{Solution.} By Kronecker--Weber, every abelian extension of $\mathbb{Q}$ is contained in $\mathbb{Q}(\zeta_n)$ for some $n$. We need $\sqrt{5}\in\mathbb{Q}(\zeta_n)$. The discriminant of $\mathbb{Q}(\sqrt{5})$ is $5$, and by class field theory, $\mathbb{Q}(\sqrt{5})\subset\mathbb{Q}(\zeta_5)$ since the quadratic subfield of $\mathbb{Q}(\zeta_5)$ is $\mathbb{Q}(\sqrt{5})$ (the unique quadratic subfield of $\mathbb{Q}(\zeta_p)$ for an odd prime $p$ is $\mathbb{Q}(\sqrt{(-1)^{(p-1)/2}p})$, which for $p=5$ gives $\mathbb{Q}(\sqrt{5})$). The smallest cyclotomic field is $\mathbb{Q}(\zeta_5)$, of degree 4 over $\mathbb{Q}$.

\item \textbf{Exercise 3.} Determine the Frobenius element $\left(\frac{\mathbb{Q}(\sqrt{5})/\mathbb{Q}}{3}\right)$ for the prime $p=3$ in the extension $\mathbb{Q}(\sqrt{5})/\mathbb{Q}$.

\emph{Solution.} We determine how $3$ splits in $\mathbb{Q}(\sqrt{5})$ by computing the Legendre symbol $\left(\frac{5}{3}\right)$. Since $5\equiv 2\pmod{3}$ and $\left(\frac{2}{3}\right)=-1$, the prime $3$ is inert in $\mathbb{Q}(\sqrt{5})$. The Frobenius element is therefore the nontrivial element of $\operatorname{Gal}(\mathbb{Q}(\sqrt{5})/\mathbb{Q})\cong\mathbb{Z}/2\mathbb{Z}$, i.e., the conjugation $\sqrt{5}\mapsto -\sqrt{5}$.

\item \textbf{Exercise 4.} For the extension $\mathbb{Q}(\zeta_7)/\mathbb{Q}$, find the decomposition of the prime $p=2$ by computing the order of $2$ modulo $7$.

\emph{Solution.} The prime $p=2$ is unramified in $\mathbb{Q}(\zeta_7)/\mathbb{Q}$. Its Frobenius element is determined by $2\bmod 7$, and the order of the Frobenius in $\operatorname{Gal}(\mathbb{Q}(\zeta_7)/\mathbb{Q})\cong(\mathbb{Z}/7\mathbb{Z})^\times$ equals the order of $2$ in $(\mathbb{Z}/7\mathbb{Z})^\times$. Compute: $2^1=2$, $2^2=4$, $2^3=8\equiv 1\pmod 7$. So the order is 3. The prime $2$ splits into $[L:\mathbb{Q}]/3=6/3=2$ primes, each with inertia degree $3$.

\item \textbf{Exercise 5.} Compute the class number of $\mathbb{Q}(\sqrt{-1})=\mathbb{Q}(i)$ by showing every ideal in $\mathbb{Z}[i]$ is principal.

\emph{Solution.} The ring $\mathbb{Z}[i]$ is a Euclidean domain with norm $N(a+bi)=a^2+b^2$: for any $\alpha,\beta\in\mathbb{Z}[i]$ with $\beta\ne 0$, there exist $q,r\in\mathbb{Z}[i]$ with $\alpha=q\beta+r$ and $N(r)<N(\beta)$. This follows from rounding $\alpha/\beta$ to the nearest Gaussian integer. Every Euclidean domain is a PID, and every PID has trivial class group. Therefore $h(\mathbb{Q}(i))=1$.

\end{enumerate}
\section*{8. References}


\begin{thebibliography}{9}
\bibitem{class_field_theory_ref1} J. Neukirch, \emph{Class Field Theory}, Springer, 1986.
\bibitem{class_field_theory_ref2} J. W. S. Cassels and A. Fröhlich, eds., \emph{Algebraic Number Theory}, Academic Press, 1967.
\bibitem{class_field_theory_ref3} J. Neukirch, A. Schmidt, and K. Wingberg, \emph{Cohomology of Number Fields}, Springer, 2nd ed., 2008.
\bibitem{class_field_theory_ref4} S. Lang, \emph{Algebraic Number Theory}, Springer, 2nd ed., 1994.
\bibitem{class_field_theory_ref5} J. W. S. Cassels and A. Fröhlich, \emph{Algebraic Number Theory}, Springer reprint.
\bibitem{class_field_theory_ref6} J. Neukirch, \emph{Algebraic Number Theory}, Springer, 1999.
\bibitem{class_field_theory_ref7} R. Greenberg, "Iwasawa theory for elliptic curves," in \emph{Arithmetic Theory of Elliptic Curves}, Springer.
\bibitem{class_field_theory_ref8} A. Grothendieck, "Brief an G. Faltings," in discussions of anabelian geometry and arithmetic fundamental groups.
\end{thebibliography}
